\subsubsection{CU-04: Hacer Preguntas sobre un Producto}
Este caso de uso describe el proceso mediante el cual un usuario podrá
formular preguntas sobre un producto específico, con el objetivo de resolver
dudas que solo se aprecian con el uso real y cotidiano del mismo antes de
realizar una compra.

\begin{table}[H]
\centering
\begin{tabular}{p{3.5cm} p{10cm} p{1cm}}\toprule \textbf{Caso de Uso} &
\textit{Hacer Preguntas sobre un Producto} & \textit{04} \\

\midrule
\textbf{Actores}       & Usuario             \\[6pt]
\textbf{Tipo}          & Primario            \\[6pt]
\textbf{Precondición}  & El usuario debe haber
escaneado previamente un producto o accedido a su página de información.  \\[6pt]
\textbf{Postcondición} & El sistema publica la pregunta en el apartado de
preguntas del producto y notifica a los usuarios que han reseñado dicho
producto.  \\[6pt]

\midrule
\multicolumn{3}{l}{\textbf{Propósito}} \\*[2pt]
\multicolumn{3}{p{14.5cm}}{Permitir que el usuario resuelva dudas específicas
sobre un producto a partir de la experiencia de otros compradores.} \\[6pt]

\midrule
\multicolumn{3}{l}{\textbf{Resumen}} \\*[2pt]
\multicolumn{3}{p{14.5cm}}{El usuario, al consultar la información de un
producto, puede acceder al apartado de dudas y redactar una nueva consulta
sobre funcionalidades que no fueron mencionadas en el apartado de experiencias o
de información técnica del producto. El usuario debera ingresar nuevamente a la 
página del producto para verificar que su pregunta ha sido respondida.} \\[6pt]
\bottomrule
\end{tabular}
\end{table}

\begin{table}[H]
\centering
\begin{tabular}{p{15.5cm}}\toprule \textbf{Descripción} \\
\midrule
\textbf{Flujo Principal} \\[2pt]
\begin{enumerate}
\item El usuario selecciona el apartado \textit{Dudas}.
\item El sistema muestra las preguntas existentes del producto junto con sus
      respuestas.
\item El usuario selecciona la opción \textit{Hacer una pregunta}.
\item El sistema muestra un campo de texto para redactar la pregunta.
\item El usuario escribe su pregunta y selecciona \textit{Publicar}.
\item El sistema valida que la pregunta cumpla con las reglas de contenido y
      longitud permitidas.
\item El sistema publica la pregunta en el apartado del producto.
\end{enumerate} \\
\midrule
\textbf{Flujo Alterno} \\[2pt]
\begin{enumerate}
\item[6a.] Si el usuario intenta publicar una pregunta vacía, el sistema
    muestra un mensaje indicando que el campo es obligatorio.
\item[7a.] Si la pregunta excede la longitud máxima permitida, el sistema
    muestra un mensaje indicando el límite de caracteres.
\item[7b.] Si la pregunta contiene contenido inapropiado u ofensivo, el
    sistema la rechaza y muestra un mensaje explicando el motivo.
\item[8a.] Si no hay conexión a internet, el sistema muestra un mensaje de error
    indicando que se requiere conexión para publicar la pregunta.
\end{enumerate}
\\
\bottomrule
\end{tabular}
\end{table}